% !TeX root = ../../Book3.tex
\section{一般双曲 Riemann 曲面}
\label{b3:sec:1304}

一般双曲曲面不必有一个单值圆盘坐标，但每个小片都可以由圆盘覆盖提供坐标。甲板不变性将长度元下降，路径提升则把距离、压缩与完备性传下来。

\subsection{从万有覆盖拉回}
\label{b3:subsec:130401}

\begin{theorem}[深层覆盖模型输入]\label{b3:thm:13-cover-model-input}
每张 Riemann 曲面的全纯万有覆盖双全纯等价于球面、平面、圆盘之一。特别地，补集至少含两个点的平面区域具有圆盘万有覆盖。
\end{theorem}
\begin{proofsketch}
主线在这里独立采用统一化输入，不要求先读选读第12章。正式来源为 Ahlfors《Conformal Invariants》\cite{Ahlfors2010Conformal}定理10-3及 \S10-4--10-6（142--151页）。平面域特例由自由甲板群分类推出：球面覆盖使底面紧，不可能；平面甲板群只能为离散平移群，底面是平面、去心平面或紧环面；平面单叶映射为仿射，去心平面单叶映射在两端由 Casorati--Weierstrass 与局部逆排除本性奇点，延为一次有理映射，所以其平面像只少一个点。补集至少两点排除所有这些例外。一般统一化的势论证明作为深层来源输入，后面的度量结论逐项完整证明。
\end{proofsketch}
称圆盘覆盖的曲面为双曲曲面。给 \(p:\mathbb D\rightarrow X\)，在底坐标中定义 \(\lambda_X(p(w))=2/((1-|w|^2)|p'(w)|)\)。若两提升为 \(w,Tw\)，由 \(p\circ T=p\) 与圆盘密度不变性得相同数值，故良定义。换底坐标时按导数模变换，所以给内在光滑长度元，局部曲率为 \(-1\)。

\subsection{双曲区域}
\label{b3:subsec:130402}

对去心盘用指数半平面覆盖直接得到
\[
 \lambda_{\mathbb D^*}(z)=\frac1{|z|\log(1/|z|)}.
\]
沿径向趋于0的长度为 \(\int\mathrm{d}r/(r\log(1/r))\)，发散；穿孔虽为一个欧氏孤立点，仍在无限距离。靠近外圆周密度也发散。

\ProseParagraphBreak
对 \(A(r,R)\)，令 \(L=\log(R/r)\)，将对数带以 \(e^{\pi\mathrm{i}(w-\log r)/L}\) 映到半平面，再下降得
\[
 \lambda_A(z)=\frac{\pi}{L|z|\sin(\pi\log(|z|/r)/L)}.
\]
中间圆 \(|z|=\sqrt{rR}\) 的长度为 \(2\pi^2/L=\pi/\operatorname{mod}A\)。任意绕孔一次的闭路长度不小于此数，因为正弦不超过1且总辐角变化至少 \(2\pi\)；中间圆达到下界。

\begin{proposition}[圆环的最短核心曲线]\label{b3:prop:13-annulus-core}
设圆环的模量为 \(m>0\)。在绕孔一次的闭曲线中，双曲长度最小值为 \(\pi/m\)，由几何平均半径 \(\sqrt{rR}\) 的圆达到。绕孔 \(n\) 次的长度至少为 \(|n|\pi/m\)。
\end{proposition}
\begin{proof}
写 \(L=\log(R/r)=2\pi m\)、\(t=\log(|z|/r)\)，密度公式给
\[
 \lambda_A(z)|\mathrm{d}z|=\frac\pi{L\sin(\pi t/L)}\frac{|\mathrm{d}z|}{|z|}.
\]
沿闭路连续选辐角 \(\theta\)，\(|\mathrm{d}z|/|z|\ge|\mathrm{d}\theta|\)，且 \(\sin(\pi t/L)\le1\)。因此长度至少为 \((\pi/L)|\theta(1)-\theta(0)|=2\pi^2|n|/L\)。在 \(t=L/2\) 的圆上，正向绕行没有角方向往返，两项不等式均等号，所以达到一次绕行的最小值。多重绕行也给相应等号。
\end{proof}
模量大意味着核心闭曲线短，而连接两端的极值长度大。两项结论衡量不同曲线族，彼此并不矛盾。若内半径趋于0，模量趋于无穷，核心长度趋于0；极限去心盘没有实际达到长度0的闭曲线，短圈越来越靠近理想端部。

\subsection{共形不变性}
\label{b3:subsec:130403}

\begin{theorem}\label{b3:thm:13-hyperbolic-schwarz-pick}
双曲曲面间全纯映射 \(f:X\rightarrow Y\) 满足 \(\lambda_Y(f(z))|f'(z)|\le\lambda_X(z)\)，并压缩路径距离。
\end{theorem}
\begin{proof}
复合 \(f\circ p_X\) 从单连通圆盘出发，可用路径提升与同伦提升取全纯提升 \(F:\mathbb D\rightarrow\mathbb D\)，使 \(p_Y\circ F=f\circ p_X\)。提升在小片由全纯层逆给出，故全纯。对 \(F\) 用圆盘 Schwarz--Pick，再以覆盖导数链式法则和密度定义消去两覆盖因子，得到所述式。沿曲线积分并取下确界给距离压缩。
\end{proof}
双全纯映射及其逆各给一个不等式，因此必须等号；一般覆盖映射也是局部等距，但不同底点的最短距离可以因不同提升选择而缩短。局部等距并不要求全局距离逐对相等。

\ProseParagraphBreak
全纯映射使长度压缩时，并没有保证不同点的像不同。\(z\mapsto z^m\) 在圆盘中是压缩映射，0处导数为0；两个不同点也可能映到同一点。若导数在某一点达到 Schwarz--Pick 等号，提升后的圆盘映射才被迫是自同构；下降到曲面以后得到的是局部等距覆盖，仍可能多层。对圆盘自映射它就是自同构，对圆环间 \(z^m\) 则恰是覆盖，显示目标拓扑影响等号结论的最终表述。

\ProseParagraphBreak
把覆盖本身看作全纯映射，密度等号处处成立，这是密度下降定义的直接后果。覆盖可以将源中的不同点识别，底面两点的最短距离是所有甲板移动后的提升距离下确界，所以不要求任意固定两条提升都保持端点距离。局部长度保持与全局距离保持有区别，后者只对选好最短相对位置的提升成立。

\subsection{单调性}
\label{b3:subsec:130404}

\begin{corollary}\label{b3:cor:13-density-monotonicity}
若双曲平面区域 \(U\subset V\)，则 \(\lambda_V\le\lambda_U\)。若 \(U\ne V\)，在每个内点都严格。
\end{corollary}
\begin{proof}
对包含映射使用前一定理。若某点取等，其圆盘提升由 Schwarz--Pick 刚性为自同构。共轭后源甲板群包含于目标甲板群，底映射是相应群商之间的覆盖，特别满射；但包含映射满射只可能 \(U=V\)。故真包含没有等号点。
\end{proof}
去心盘包含于圆盘，直接公式也给 \(\lambda_{\mathbb D^*}>\lambda_{\mathbb D}\)。删点使全纯映射的活动空间更小，所以局部内在密度增大，而不是减小。

\begin{proposition}[凸域中的边界距离估计]\label{b3:prop:13-convex-boundary-distance}
设 \(U\subsetneq\mathbb C\) 为凸区域，\(d(z)=\operatorname{dist}(z,\partial U)\)。则
\[
 \frac1{d(z)}\le\lambda_U(z)\le\frac2{d(z)}.
\]
\end{proposition}
\begin{proof}
以 \(z\) 为中心、半径 \(d(z)\) 的盘包含于 \(U\)，密度单调性给上界。取最近边界点 \(q\)。凸性使经过 \(q\)、垂直于 \(z-q\) 的直线为支持线，\(U\) 落在含 \(z\) 的开半平面 \(H\) 中。可从闭凸集在 \(q\) 的支持直线存在性及内切圆的切线唯一性看出这一点；这里使用实平面的基本凸性分离事实。半平面密度在 \(z\) 为 \(1/d(z)\)，再用单调性得下界。
\end{proof}
半平面达到下界，圆盘中心达到上界，故两个常数在凸域类中不能整体改进。对于任意非凸域，上界仍由内盘包含成立，下界却不能不经新证明沿用。去心盘靠近穿孔时 \(\lambda\sim1/(|z|\log(1/|z|))\)，Euclid 边界距离为 \(|z|\)，已小于 \(1/d(z)\)，显示凸性确实参与下界。

\ProseParagraphBreak
下界也解释凸边界的无限距离。边界距离函数是1-Lipschitz，沿任意路径 \(|\mathrm{d}d|\le|\mathrm{d}z|\)，所以双曲长度至少为 \(\int|\mathrm{d}d|/d\)。若路径从固定正距离走到边界距离0，该积分至少为相应对数差并趋于无穷。这里无需假定路径径向，也无需边界可微；凸性及支持半平面已经提供整个路径上的下界。

\subsection{完备性}
\label{b3:subsec:130405}

\begin{theorem}\label{b3:thm:13-hyperbolic-complete}
双曲曲面的 Poincaré 路径度量完备。任意两点有最短测地线，且对选定提升 \(u,v\) 有
\[
 \mathrm d_X(pu,pv)=\inf_{T\in\operatorname{Deck}(p)}\mathrm d_{\mathbb D}(u,Tv).
\]
\end{theorem}
\begin{proof}
任意底路径从 \(u\) 提升，终点是某个 \(Tv\)，长度保持，所以底距离至少右端；任意圆盘最短路投影给底路径，反向不等式成立。纤维是闭离散集，圆盘闭有界球紧，故在给定有限距离内只有有限个 \(Tv\)；下确界由其中一点达到，最短路投影得到底测地线。

设底空间有 Cauchy 列，抽取子列 \(x_n\) 满足相邻距离小于 \(2^{-n}\)，选连接路径长度小于 \(2^{1-n}\)。从选定提升开始逐次提升这些路径，得到圆盘点 \(u_n\)。尾部连接长度和趋于0，使 \(u_n\) 在圆盘为 Cauchy，完备性给极限 \(u\)。投影后 \(x_n\rightarrow p(u)\)，原 Cauchy 列也趋于同一点。因此底度量完备。
\end{proof}
在环域中不同绕行类可能给不同测地线，所以唯一性不再像圆盘那样自动成立。去心盘的绕孔闭路长度可趋于0，却没有长度0的闭测地线；完备性不能误解为每个非平凡闭路类都有长度极小代表。

\ProseParagraphBreak
覆盖距离公式还能给局部没有自重叠的半径。对 \(x=p(w)\)，令
\[
 r_x=\tfrac12\inf_{T\ne1}d_{\mathbb D}(w,Tw),
\]
若甲板群平凡便取无穷。自由离散作用和闭球紧使这个下确界为正：否则非平凡轨道点将趋于 \(w\)，违反覆盖均匀邻域的分离。小于 \(r_x\) 的圆盘球投影单射；若两点被 \(T\) 识别，三角不等式会使 \(d(w,Tw)<2r_x\)，矛盾。这给曲面上围绕 \(x\) 的一张无重叠等距坐标图。

\ProseParagraphBreak
对于有限圆环，核心附近的这种半径由核心长度控制；对去心盘，向穿孔端行进时绕行圈的长度趋于0，所以局部单射半径也可能趋于0。曲面仍完整，因为通向穿孔的径向长度 \(\int\mathrm{d}r/[r\log(1/r)]\) 发散。越来越短的绕行闭路与越来越远的端部可以同时存在，完整度量并没有要求全曲面的局部单射半径有统一正下界。

\begin{proposition}[商曲面的闭球与最短路]\label{b3:prop:13-quotient-geodesics}
圆盘型曲面的有限半径闭球紧，每两点间都有长度等于距离的路径；最短路径不必唯一。
\end{proposition}
\begin{proof}
固定 \(x=p(w)\)。覆盖距离公式及离散纤维说明，对任意 \(y\)，可在其纤维中找到 \(v\) 使 \(d_X(x,y)=d_{\mathbb D}(w,v)\)：若距离下确界有限，取一个有界距离序列，圆盘闭球紧而纤维闭离散，故其中只有有限个点，下确界达到。因此半径 \(R\) 的闭球恰是 \(p(\overline B_{\mathbb D}(w,R))\)，是紧集。把 \(w\) 到最佳提升 \(v\) 的圆盘测地段投影，其长度保持并等于底面距离，故是最短路。
\end{proof}
有限圆环的核心圆上两个相对点有两条最短路径，分别沿核心走半圈。任意连接路径的连续辐角变化为 \(\pi+2\pi n\)，绝对值至少 \(\pi\)；核心长度证明中的角积分下界说明半圈达到最小值。圆盘自身的最短路唯一，商以后却因不同甲板相对位置取得同样最小值而不唯一。局部测地线仍由圆盘模型决定，全局唯一性不随覆盖自动下降。
